\documentclass[11pt,a4paper]{article}
\usepackage[margin=25mm]{geometry}
\usepackage{fontspec}
\tracinglostchars=3
\setmainfont{Latin Modern Roman}
\setsansfont{Latin Modern Sans}
\setmonofont{DejaVu Sans Mono}[Scale=0.78]
\usepackage{amsmath,amssymb,amsthm,mathtools}
\usepackage{longtable,booktabs,array,calc,enumitem}
\usepackage{xcolor}
\usepackage[unicode,colorlinks=true,linkcolor=blue!45!black,urlcolor=blue!55!black]{hyperref}
\usepackage{xurl}
\usepackage{newunicodechar}
\newunicodechar{ℓ}{\ensuremath{\ell}}
\newunicodechar{π}{\ensuremath{\pi}}
\newunicodechar{χ}{\ensuremath{\chi}}
\newunicodechar{φ}{\ensuremath{\varphi}}
\newunicodechar{ζ}{\ensuremath{\zeta}}
\newunicodechar{′}{\ensuremath{\prime}}
\newunicodechar{α}{\ensuremath{\alpha}}
\newunicodechar{β}{\ensuremath{\beta}}
\newunicodechar{γ}{\ensuremath{\gamma}}
\newunicodechar{δ}{\ensuremath{\delta}}
\newunicodechar{ε}{\ensuremath{\varepsilon}}
\newunicodechar{θ}{\ensuremath{\theta}}
\newunicodechar{λ}{\ensuremath{\lambda}}
\newunicodechar{μ}{\ensuremath{\mu}}
\newunicodechar{ν}{\ensuremath{\nu}}
\newunicodechar{ξ}{\ensuremath{\xi}}
\newunicodechar{ρ}{\ensuremath{\rho}}
\newunicodechar{σ}{\ensuremath{\sigma}}
\newunicodechar{τ}{\ensuremath{\tau}}
\newunicodechar{ψ}{\ensuremath{\psi}}
\newunicodechar{ω}{\ensuremath{\omega}}
\newunicodechar{Γ}{\ensuremath{\Gamma}}
\newunicodechar{Δ}{\ensuremath{\Delta}}
\newunicodechar{Θ}{\ensuremath{\Theta}}
\newunicodechar{Λ}{\ensuremath{\Lambda}}
\newunicodechar{Ξ}{\ensuremath{\Xi}}
\newunicodechar{Π}{\ensuremath{\Pi}}
\newunicodechar{Σ}{\ensuremath{\Sigma}}
\newunicodechar{Φ}{\ensuremath{\Phi}}
\newunicodechar{Ψ}{\ensuremath{\Psi}}
\newunicodechar{Ω}{\ensuremath{\Omega}}
\newunicodechar{₀}{\ensuremath{_0}}
\newunicodechar{₁}{\ensuremath{_1}}
\newunicodechar{₂}{\ensuremath{_2}}
\newunicodechar{∞}{\ensuremath{\infty}}
\newunicodechar{≤}{\ensuremath{\leq}}
\newunicodechar{≥}{\ensuremath{\geq}}
\newunicodechar{∈}{\ensuremath{\in}}
\newunicodechar{⊗}{\ensuremath{\otimes}}
\newunicodechar{≅}{\ensuremath{\cong}}
\newunicodechar{≃}{\ensuremath{\simeq}}
\newunicodechar{∼}{\ensuremath{\sim}}
\newunicodechar{≠}{\ensuremath{\ne}}
\newunicodechar{→}{\ensuremath{\to}}
\newunicodechar{↔}{\ensuremath{\leftrightarrow}}
\newunicodechar{⊂}{\ensuremath{\subset}}
\newunicodechar{∩}{\ensuremath{\cap}}
\newunicodechar{∪}{\ensuremath{\cup}}
\usepackage{microtype}
\setlength{\emergencystretch}{3em}
\setlength{\parskip}{0.35em}
\setlength{\parindent}{0pt}
\setcounter{tocdepth}{2}
\providecommand{\tightlist}{\setlength{\itemsep}{0pt}\setlength{\parskip}{0pt}}
\newtheorem{theorem}{Theorem}[section]
\newtheorem{proposition}[theorem]{Proposition}
\newtheorem{lemma}[theorem]{Lemma}
\newtheorem{corollary}[theorem]{Corollary}
\theoremstyle{definition}\newtheorem{definition}[theorem]{Definition}
\theoremstyle{remark}\newtheorem{remark}[theorem]{Remark}
\title{ALPOGE\_BRIDGES\\The Riemann hypothesis and Dirichlet GRH\\in Grothendieck's programme}
\author{Research synthesis and proof checks in the present run}
\date{21 November 2026\quad\textnormal{(Europe/Berlin)}}
\begin{document}
\maketitle
\begin{abstract}
We adopt the supplied Alp\"oge RH and Dirichlet-GRH declarations as
theorems and determine their precise analytic and motivic consequences.
The main work concerns questions and unfinished statements of Grothendieck
and Deligne, using original archive sheets and exact literature locators.
Full proofs are supplied for the deductions and counterexamples established
in this run. Existing results are cited with their operative hypotheses.
The arithmetic consequences include shifted critical lines for abelian
Artin--Tate objects, effective finite Frobenius recognition, the corresponding
Borger realizations, explicit classical consequences, and trace positivity.
Separate cycle-theoretic, Hodge-theoretic and categorical arguments settle
several archive statements without using RH. Counterexamples locate the
information lost by passing to semisimple Euler factors or numerical traces.
Unresolved questions retain explicit gates and attempted approaches.
This working manuscript records ongoing work; its inclusion of a gate
is not a claim that the corresponding problem is solved.
\end{abstract}

\section*{Reading the claims and the check record}
The mathematical input is Alp\"oge (2026), \emph{Remark}, together with
the five named declarations in \texttt{ALPOGE\_Main.lean} and their
targets in \texttt{ALPOGE\_Statements.lean}. The supplied Mathlib revision
is recorded in the source-scope section. The declarations have been
inspected; the imported proof tree is not in the supplied packet, and
no fresh Lean kernel compilation is claimed. This does not weaken or
replace the input adopted by the task. In this manuscript, ``under the
adopted theorem'' describes the provenance of a consequence; it is not
an additional unproved RH hypothesis within the task's adopted theory.

An identifier denotes a finding established in this run, not a claim
of first publication. We do not claim an exhaustive novelty search.
The check table records actual full-proof readings, primary-source
statement checks, computations, and carried background separately.
Prior intake check reports are never relabelled as checks performed
here. The word ``refuted'' always names the precise statement or
proposed inference for which a counterexample is given.

The strongest arithmetic results concern primitive Dirichlet factors
and integral Tate twists. General Hecke characters over other fields,
nonabelian Artin representations, ramified mixed extensions, integral
lattices, period torsors and monodromy require the additional data
specified in their sections. A proved failure of one reconstruction
rule is not a logical-independence theorem about RH.

\section*{Ranking by the question settled}
\begin{enumerate}
\item \textbf{Original archive statements.} The cote 14 implication and
its proposed converse are addressed over an algebraically closed field
for the precisely defined generated Albanese quotient. All smooth
quasiprojective parameters are treated in every characteristic for the
cohomological implication; their integral relation quotient is zero in
characteristic zero and is proved to be primary torsion at the
characteristic in positive characteristic. The finite-field rational-point reading
has a separate counterexample. Cote 15's abelian--Artin motivic sector is
now treated using the cited 2026 all-CM Hodge theorem and its exact Tate
and standard-conjecture consequences; arbitrary motives and rational
fibre functors retain their distinct gates. Cote 18's real-Frobenius,
weight-shift and strong-period formulations, and cote 67's capping and
involution claims, receive explicit proofs or counterexamples.
The canonical capping map has an exact real-analytic image but fails
the natural holomorphic formulation. Cote 123's dimension, rational
singularity, nilpotent uniqueness and differential-rank questions have
exact classical answers and checked short deductions. Its strict
nonsimply-laced subregular rank prediction fails: the rank is always
one less than the semisimple rank in every type. The group quotient has a complete
central-character formula; type A has a closed minimal-polynomial
formula. The apposition quadruples are classified completely in type A.
Cote 66's named affine torsion-coefficient comparison cases are
affirmative; its universal degree-zero test is exactly henselianity.
\item \textbf{Deligne and the geometric gate.} The actual problem
Weil II 1.8.15 is answered by the precise admissible-Hodge-theory
theorems. Singular-fibre weights and nearby-cycle weights are computed
separately. The independently checked finite exterior-power criterion reduces the weight--monodromy gate to specified
inertia-invariant trace bounds. The universal stable curve has an explicit
geometric perverse decomposition, and its actual exit functor has the
proved right-Kan direct-image description. The linear theorem of cote
134-8 is established for all countable scalar rings, more generally
under a finite scalar pure-injective bound, or under either of two
finite-generation conditions, with the specified exact adjunction,
scalar compatibility and normalization retained. The unrestricted
large-ring case keeps its explicit counit gate.
\item \textbf{New effective arithmetic access.} AQ-FINITE gives a
finite recognition bound for semisimple abelian Artin--Tate objects;
BR-AT-01 and CC-F01/F01b give the exact shifted analytic factors.
The effective Siegel deduction, abelian-field Chebotarev and prime-search
bounds discharge their original conditional hypotheses.
\item \textbf{Criteria and failed routes.} Trace-class positivity,
moment inequalities, Li and approximation criteria are distinguished
from independent geometric constructions. Explicit mixed-extension,
monodromy, cross-ratio, divided-power and spectral counterexamples
identify the missing information. The 2026 rational Deligne--Drinfeld
theorem is an external result, separate from profinite GT surjectivity.
\end{enumerate}
\tableofcontents
\clearpage
\part{Checked findings and exactly scoped known results}
